\documentclass[12pt,a4paper]{article}
\IfFileExists{pt-common.sty}{\usepackage{pt-common}}{\usepackage{../shared/pt-common}}
\usepackage[margin=2.5cm]{geometry}
\usepackage{setspace}
\usepackage{tikz-3dplot}
\onehalfspacing

\title{\textbf{The Theory of Persistence~II:\\Information Geometry and Holonomy}}
\author{Yan Senez\\\texttt{yan.senez@protonmail.com}}
\date{April 2026}

\begin{document}
\maketitle

\begin{abstract}
This is the second of five articles presenting the mathematical
foundations of the Theory of Persistence (PT).  Building on the sieve
dynamics, uniqueness, conservation laws, and Gap Fundamental Theorem
established in the companion article~\cite{companion_M1}, we develop
two pillars of the theory: the canonical information geometry of the
sieve state space, and the holonomy structure that emerges from the
sieve's transfer matrices.

\textbf{Information geometry.}
The KL divergence $\DKL$ is the unique $f$-divergence compatible with
the CRT product structure of the sieve state space (Theorem~G1, via
Shore--Johnson reduction~\cite{ShoreJohnson1980}, with 7/7 elimination
tests; an autonomous CRT proof via Csisz\'ar's
characterisation~\cite{Csiszar1967} is also given).  The Fisher metric
$\Fisher$ emerges as the second variation of $\DKL$ (Theorem~G2,
standard) and is the unique monotone Riemannian metric on $\Delta_2$
(\v{C}encov reduction~\cite{Cencov1982}, Theorem~G3, 7/7 elimination
tests).  We exhibit the Fisher metric on binary and ternary simplices,
show it is forced by the optimal expander structure of the mod-3
transition graph, and connect it to the addition/multiplication duality
of the integers.

\textbf{Holonomy.}
The holonomy angle $\theta_p$ satisfies the exact algebraic identity
$\sinsp{p} = \delta_p(2 - \delta_p)$, derived by three independent
routes (geometric, spectral, information-geometric).  The anomalous
dimensions $\gamma_p = -d\ln\sinsp{p}/d\ln\mu$ are strictly decreasing
in $p$; the active prime criterion $\gamma_p > s = 1/2$ selects exactly
$\{3, 5, 7\}$ at the fixed point $\mustar = 15$, yielding $N_c = 3$.
The Casimir coefficient $C_F = (N_c^2 - 1)/(2N_c) = 4/3$ is a
consequence, not an input.  An independent algebraic check gives
$(N_c - 3)(N_c + 1) = 0$.  Ghost primes ($p \geq 11$) contribute
perturbatively as vacuum polarization corrections.

All results follow from $s = 1/2$ and the sieve structure established
in~\cite{companion_M1}.  No physics is invoked; no parameter is fitted.
\end{abstract}

\tableofcontents
\newpage

% ============================================================
%  Introduction
% ============================================================
\section{Introduction}
\label{sec:intro}

This article is the second instalment (M2) in a five-part series
developing the mathematical foundations of the Theory of Persistence
(PT).  The series is organised as follows:
\begin{description}[nosep]
\item[M1~\cite{companion_M1}:] The sieve as a dynamical field,
  uniqueness of Eratosthenes, conservation laws, and the Gap
  Fundamental Theorem.
\item[M2 (this article):] Information geometry, holonomy, and the
  internal angle scale.
\item[M3~\cite{companion_M3}:] Convergence, exact recurrences, and
  the fixed point $\mustar = 15$.
\item[M4~\cite{companion_M4}:] Extended mathematical structures
  (complex mechanics, spectral graphs, symmetry).
\item[M5~\cite{companion_M5}:] The bridge from arithmetic to physics.
\end{description}
\noindent
A separate companion~\cite{companion_physics} develops the detailed
physical predictions (Feynman programme, 43 observables);
M5 provides the structural bridge, while~\cite{companion_physics}
provides the quantitative validation.

The companion article~\cite{companion_M1} established the foundational
layer.  We recall the essential results:
\begin{itemize}[nosep]
\item \textbf{T1 (forbidden transitions):} the transfer matrix
  $T_3 = \mathrm{antidiag}(1,1)$ has a unique stationary distribution
  at $(1/2, 1/2)$, giving $s = 1/2$ with no free parameter.
\item \textbf{T6 (uniqueness):} three independent proofs that
  Eratosthenes is the unique self-consistent sieve.
\item \textbf{L0 (conservation):} the maximum-entropy gap distribution
  is geometric; the conservation exponent is $\alpha_{\rm cons} = s^2
  = 1/4$.
\item \textbf{GFT:} the exact identity $\log_2 m = \DKL + H$
  decomposes total capacity into structure and noise.
\item \textbf{CRT gauge connection:} the Chinese Remainder Theorem
  endows the sieve state space $\Delta_m$ with a product structure
  $\Delta_m \cong \prod_{p|m} \Delta_p$.
\end{itemize}

The present article develops two themes.

\textbf{Information geometry (Section~\ref{sec:geometry}).}
We prove that the KL divergence $\DKL$ and the Fisher metric $\Fisher$
are the \emph{unique} divergence and metric compatible with the sieve's
symmetries.  The proofs are reductions to classical
theorems---Shore--Johnson~\cite{ShoreJohnson1980} for $\DKL$ and
\v{C}encov~\cite{Cencov1982} for $\Fisher$---after verifying that the
sieve state space satisfies their premises.  An autonomous proof of G1
via Csisz\'ar's characterisation of additive $f$-divergences is also
given.  This section is entirely mathematical and invokes no physics.

\textbf{Holonomy (Section~\ref{sec:holonomy}).}
The sieve's transfer matrix at each prime $p$ generates a rotation
angle $\theta_p$ on the residue space $\mathbb{Z}/p\mathbb{Z}$.  The
holonomy identity $\sinsp{p} = \delta_p(2 - \delta_p)$ is an algebraic
tautology, verified by three independent routes.  The anomalous
dimensions $\gamma_p$ are strictly decreasing, and the criterion
$\gamma_p > s = 1/2$ selects exactly three active primes
$\{3, 5, 7\}$ at $\mustar = 15$, giving $N_c = 3$ and
$C_F = 4/3$.  The physical identification of these arithmetic
quantities with gauge couplings and QCD colours is a bridge
claim deferred to~\cite{companion_M5}.

\medskip\noindent\textbf{Forward pointers.}
The companion article~\cite{companion_M3} proves convergence of the
sieve recurrences and establishes $\mustar = 15$ as the unique
non-trivial fixed point.  The extended mathematical structures
(complex mechanics, spectral graphs, and symmetry extensions) are
developed in~\cite{companion_M4}.  The bridge from arithmetic to
physics---where $\sinsp{p}$ becomes a coupling constant and $N_c = 3$
becomes the number of QCD colours---is the subject
of~\cite{companion_M5}.


% ============================================================
%  S6. CANONICAL INFORMATION GEOMETRY
% ============================================================

\section{Canonical Information Geometry}
\label{sec:geometry}

\epigraph{Information geometry is a method of exploring the world of
information by means of modern geometry.}{Shun-ichi Amari, \textit{Information Geometry and Its Applications}, Springer (2016), Preface}

% ---- Status Box (R1) ----------------------------------------
\begin{statusbox}
\textbf{GOAL:}
  Prove that the two canonical divergences on the sieve state space---
  the KL divergence $\DKL$ and the Fisher metric $\Fisher$---are the
  \emph{unique} choices compatible with the sieve's symmetries.

\textbf{INPUTS:}
  CRT, gauge connection, state space $\Delta_m$ (\cite{companion_M1}).
  GFT, $\DKL$ already used (\cite{companion_M1}).

\textbf{MAIN RESULTS:}
  G1: $\DKL$ is the unique $f$-divergence invariant under CRT (Shore--Johnson
  reduction, 7/7).
  G2: Second variation of $\DKL$ = Fisher metric (standard, Amari 1985).
  G3: Fisher metric is the unique monotone Riemannian metric on $\Delta_2$
  (\v Cencov reduction, 7/7).

\textbf{STATUS:}
  \THMTAG~(reduction: G1 via Shore--Johnson, G3 via \v Cencov).
  This chapter does NOT prove Shore--Johnson or \v Cencov autonomously.
  It verifies that the sieve state space satisfies their premises.

\textbf{NOT PROVED:}
  Shore--Johnson (1980) and \v Cencov (1982): imported as named classical
  theorems. Physical use of Fisher metric (the companion article~\cite{companion_physics}).
\end{statusbox}

% ============================================================
\subsection{Intuition: why these two objects?}
\label{sec:ch5_intuition}

The Gap Fundamental Theorem (see~\cite{companion_M1}) uses $\DKL$ as the measure of
sieve structure. The companion article~\cite{companion_physics} will use the Fisher metric to derive general
relativity. But why $\DKL$ and Fisher? Among the infinitely many
possible divergences and metrics on probability distributions, what
singles these two out?

The answer, for $\DKL$: it is the unique $f$-divergence that is
additive under the CRT factorization that the sieve enjoys. Any other
$f$-divergence (Hellinger, $\chi^2$, total variation) fails this
invariance (verified: 7/7 tests).

The answer, for Fisher: it is the unique Riemannian metric on probability
simplices that contracts under Markov maps. Since the sieve's transfer
matrix $T_m$ is a Markov map, any information-geometric distance must
contract under it. Fisher is the unique metric with this property.

Both results are reductions to classical theorems (Shore--Johnson~\cite{ShoreJohnson1980},
\v{C}encov~\cite{Cencov1982}). PT's contribution is verifying that the sieve state
space satisfies the premises.

% ============================================================
\subsection{The sieve state space}
\label{sec:ch5_state_space}

\begin{definition}[Probability simplex $\Delta_m$]
\label{def:simplex}
The \emph{sieve state space} at modulus $m$ is the probability simplex:
\[
  \Delta_m = \Bigl\{ P = (p_0, \ldots, p_{m-1}) \in \mathbb{R}^m
             : p_r \geq 0,\; \sum_r p_r = 1 \Bigr\}.
\]
The sieve distribution $P_m \in \Delta_m$ is the vector of prime
residue frequencies modulo $m$.
\end{definition}

The sieve acts on $\Delta_m$ via the CRT factorization:
for $m = q_1 \cdots q_k$ (square-free), the CRT isomorphism gives
\[
  \Delta_m \cong \Delta_{q_1} \times \cdots \times \Delta_{q_k},
\]
so the state space is a Cartesian product of smaller simplices.
Any geometric object (divergence, metric) on $\Delta_m$ must respect
this product structure.
Figure~\ref{fig:fisher_simplex} illustrates the Fisher metric on the
binary simplex $\Delta_2$, and Figure~\ref{fig:fisher_3d} shows the
extension to the ternary simplex $\Delta_3$.

\begin{figure}[htbp]
\centering
\vspace{0.5cm}
\begin{tikzpicture}[scale=5.0, >=stealth]
  % Euclidean baseline
  \draw[gray!30, thick] (0,0) -- (1,0);
  \node[below, font=\small, gray] at (0,-0.02) {$0$};
  \node[below, font=\small, gray] at (1,-0.02) {$1$};
  \node[below, font=\small, gray] at (0.5,-0.02) {$\tfrac{1}{2}$};
  \node[font=\scriptsize, gray] at (0.5,-0.09) {(Euclidean)};
  % Fisher arc
  \draw[thmblue, very thick, domain=0:180, samples=80]
    plot ({0.5 + 0.45*cos(\x)}, {0.45*sin(\x)});
  \node[above, font=\small, thmblue] at (0.5, 0.53) {Fisher arc};
  % Grid lines showing metric stretching near boundaries
  \foreach \p in {0.05, 0.1, 0.2, 0.3, 0.4, 0.5, 0.6, 0.7, 0.8, 0.9, 0.95} {
    \pgfmathsetmacro{\ang}{asin(sqrt(\p))*2}
    \pgfmathsetmacro{\fx}{0.5 + 0.45*cos(\ang)}
    \pgfmathsetmacro{\fy}{0.45*sin(\ang)}
    \draw[thmblue!40, thin] (\p, 0) -- (\fx, \fy);
  }
  % Stationary point p = 1/2 — leader line below Fisher arc label
  \fill[condorange] (0.5, 0.45) circle (0.7pt);
  \draw[condorange, thin] (0.5, 0.45) -- (0.68, 0.48);
  \node[right, font=\small, condorange] at (0.68, 0.48) {$s = \tfrac{1}{2}$};
  % Stretching near boundaries
  \draw[<->, thin, predred] (0.02, -0.12) -- (0.12, -0.12);
  \node[below, font=\scriptsize, predred] at (0.07, -0.13) {$ds \to \infty$};
  \draw[<->, thin, predred] (0.88, -0.12) -- (0.98, -0.12);
  \node[below, font=\scriptsize, predred] at (0.93, -0.13) {$ds \to \infty$};
  % Metric formula
  \node[font=\normalsize] at (0.5, -0.17)
    {$ds^2 = \dfrac{dp^2}{p\,(1-p)}$};
  % axis label
  \node[below, font=\small] at (0.5, -0.28) {probability $p = \pi_1$};
\end{tikzpicture}
\caption{Fisher metric on the binary simplex $\Delta_2$.
  The Euclidean interval $[0,1]$ (gray) is ``curved'' by the Fisher
  metric $ds^2 = dp^2/[p(1-p)]$ into a half-circle (blue arc).
  Grid lines connect equal probability values: they crowd near the
  boundaries ($p \to 0$ or $p \to 1$), showing the metric divergence.
  The stationary point $s = 1/2$ (orange) sits at the geodesic midpoint.
  This is the \emph{unique} monotone metric on $\Delta_2$ (\v Cencov).}
\label{fig:fisher_simplex}
\end{figure}

\begin{figure}[htbp]
\centering
\tdplotsetmaincoords{68}{125}
\begin{tikzpicture}[tdplot_main_coords, scale=4.0, >=stealth]
  % Ternary simplex Delta_3: vertices
  \coordinate (P1) at (1,0,0);
  \coordinate (P2) at (0,1,0);
  \coordinate (P3) at (0,0,1);
  % Simplex edges
  \draw[gray!50, thick] (P1) -- (P2) -- (P3) -- cycle;
  % Fill simplex face
  \fill[thmblue!8, opacity=0.4] (P1) -- (P2) -- (P3) -- cycle;
  % Grid lines on the simplex (explicit coordinates, no calc)
  \foreach \i in {1,...,8} {
    \pgfmathsetmacro{\t}{\i/9}
    \pgfmathsetmacro{\omt}{1-\t}
    % Lines parallel to P1-P2 edge
    \draw[thmblue!30, very thin]
      ({\omt}, 0, {\t}) -- (0, {\omt}, {\t});
    % Lines parallel to P2-P3 edge
    \draw[thmblue!30, very thin]
      ({\t}, {\omt}, 0) -- ({\t}, 0, {\omt});
    % Lines parallel to P1-P3 edge
    \draw[thmblue!30, very thin]
      (0, {\t}, {\omt}) -- ({\omt}, {\t}, 0);
  }
  % Centroid = (1/3, 1/3, 1/3)
  \fill[condorange] (0.333, 0.333, 0.333) circle (0.8pt);
  \draw[condorange, thin] (0.333, 0.333, 0.333) -- (0.6, 0.6, 0.5);
  \node[right, font=\small, condorange] at (0.6, 0.6, 0.5)
    {$p_i = \tfrac{1}{3}$ ($s = \tfrac{1}{2}$)};
  % Vertex labels
  \node[below right, font=\small] at (P1) {$p_1 = 1$};
  \node[below left, font=\small] at (P2) {$p_2 = 1$};
  \node[above, font=\small] at (P3) {$p_3 = 1$};
  % Edge labels (divergence)
  \node[font=\scriptsize, predred] at (0.5, 0.5, 0) {$ds \to \infty$};
  \node[font=\scriptsize, predred] at (0, 0.5, 0.5) {$ds \to \infty$};
  \node[font=\scriptsize, predred] at (0.5, 0, 0.5) {$ds \to \infty$};
  % Metric formula — offset left to avoid overlapping triangle
  \node[font=\small, align=center, gray] at (-0.6, -0.2, 0.5)
    {$ds^2 = \sum_i \dfrac{dp_i^2}{p_i}$};
\end{tikzpicture}
\caption{Fisher metric on the ternary simplex $\Delta_3$.
  The three vertices (pure states, $p_i = 1$) are at infinite Fisher
  distance from any interior point---the metric diverges near edges.
  Grid lines show how equal Euclidean steps correspond to unequal Fisher
  lengths (crowding near boundaries). The centroid $p_i = 1/3$ (orange,
  corresponding to $s = 1/2$ in the binary projection) is the geodesic
  midpoint.}
\label{fig:fisher_3d}
\end{figure}

% ============================================================
\subsection{G1: Uniqueness of $\DKL$ via Shore--Johnson}
\label{sec:G1}

\paragraph{The Shore--Johnson theorem (1980, external import).}
Under the following axioms on a divergence $D(\cdot\|\cdot)$:
\begin{enumerate}[label=(SJ\arabic*)]
\item \textbf{Consistency}: $\arg\min_Q D(Q\|P)$ subject to constraints
  is independent of the order in which constraints are imposed.
\item \textbf{Invariance}: $D$ is invariant under permutation of labels.
\item \textbf{System independence}: for independent systems, $D$ is additive.
\item \textbf{Subset independence}: updating on a constraint that concerns
  only a subset changes only that subset's distribution.
\item \textbf{Scaling}: $D$ transforms consistently under relabeling.
\end{enumerate}
Shore and Johnson~\cite{ShoreJohnson1980} proved: the unique divergence satisfying SJ1--SJ5
and yielding the maximum entropy principle is $\DKL$.

\mTHM
\begin{theorem}[KL Divergence Uniqueness (G1)]
\label{thm:G1}
On the sieve state space $\Delta_m$ with the CRT product structure:
$\DKL$ is the unique $f$-divergence satisfying Shore--Johnson axioms
SJ1--SJ5 and invariant under CRT factorization.

\begin{proof}
\textbf{Verification of premises.} The sieve state space $\Delta_m$
with the product structure $\Delta_m \cong \prod_{p|m} \Delta_p$ satisfies:
\begin{itemize}
\item SJ3 (system independence): CRT factorization gives $\DKL(\prod P_p \| \prod Q_p)
  = \sum_p \DKL(P_p \| Q_p)$. Independence of the CRT components implies
  additivity. (Verified: 7/7 tests.)
\item SJ4 (subset independence): updating residues modulo $p$ changes only the
  $p$-component in the CRT decomposition. (Verified: 7/7 tests.)
\item SJ1, SJ2, SJ5: standard consistency/invariance properties satisfied by
  the uniform reference $U_m$ and the empirical distribution $P_m$.
\end{itemize}

\textbf{Application of Shore--Johnson.} Since the premises SJ1--SJ5 hold
on $\Delta_m$, the Shore--Johnson theorem gives: $\DKL$ is the unique
divergence satisfying these axioms. (External import: Shore--Johnson 1980,
published in IEEE Trans.\ Inf.\ Theory.)

\textbf{Elimination of alternatives.}
\begin{itemize}
\item $\chi^2$ divergence: violates SJ3 (not additive under CRT, ratio
  of squared differences not additive).
\item Hellinger distance: violates SJ3 (square root breaks additivity).
\item Total variation: violates SJ4 (updating one coordinate affects
  the $\ell_1$ norm globally).
\item R\'enyi divergences $D_\alpha$ ($\alpha \neq 1$): violate SJ1
  (inconsistency of sequential updates for $\alpha \neq 1$).
\end{itemize}
\end{proof}
\end{theorem}

\begin{remark}[Scope of G1]
G1 is a \emph{reduction} to Shore--Johnson, not an autonomous proof.
The PT contribution: verifying that $\Delta_m$ with CRT structure
satisfies SJ1--SJ5. The uniqueness conclusion is imported.
\end{remark}

\begin{theorem}[CRT-Autonomous KL Uniqueness (G1$'$)]
\label{thm:G1_autonomous}
On the sieve state space $\Delta_m$ with CRT product structure, $\DKL$
is the unique $f$-divergence satisfying additivity under independent subsystems.

\begin{proof}
\textbf{Step 1 (CRT $\Rightarrow$ product structure).}
By the Chinese Remainder Theorem (see Theorem~3.10 in~\cite{companion_M1}), for
$m = q_1 \cdots q_k$ (square-free):
$\Delta_m \cong \Delta_{q_1} \times \cdots \times \Delta_{q_k}$.
Any divergence $D$ on $\Delta_m$ must respect this factorization.

\textbf{Step 2 (Additivity characterisation).}
An $f$-divergence $D_f(P\|Q) = \sum_i q_i\,f(p_i/q_i)$ is additive on
product distributions ($P = P_1 \otimes P_2$, $Q = Q_1 \otimes Q_2$
$\Rightarrow D_f(P\|Q) = D_f(P_1\|Q_1) + D_f(P_2\|Q_2)$) if and only
if $f$ has the form $f(t) = c_1\,t\ln t + c_2(t-1)$ for constants $c_1 > 0$,
$c_2 \in \mathbb{R}$ (Csisz\'ar~\cite{Csiszar1967}).
Since $\sum_i q_i(p_i/q_i - 1) = 0$ on normalised distributions,
the $c_2$ term vanishes identically. The unique additive $f$-divergence
(up to positive rescaling) is therefore
$D_f(P\|Q) = \sum_i p_i \ln(p_i/q_i) = \DKL(P\|Q)$.

\textbf{Step 3 (Uniqueness).}
Since CRT imposes product structure at every sieve level, and since
$f(t) = t\ln t$ is the unique generator producing an additive divergence,
$\DKL$ is the unique additive $f$-divergence on $\Delta_m$.
No Shore--Johnson import is needed; the proof uses only CRT, the
definition of $f$-divergences, and Csisz\'ar's characterisation
of additive generators (elementary convex analysis).

\textbf{Elimination.}
\begin{itemize}[nosep]
\item $\chi^2$: $f(t) = (t-1)^2$. Not of the form $c_1 t\ln t + c_2(t-1)$.
  Numerical: $D_{\chi^2}(P\otimes P' \| Q\otimes Q') \neq D_{\chi^2}(P\|Q) + D_{\chi^2}(P'\|Q')$ (verified).
\item Hellinger: $f(t) = (\sqrt{t}-1)^2$. Not of the required form. Fails additivity (verified).
\item Total variation: $f(t) = |t-1|$. Not differentiable at $t=1$. Fails additivity (verified).
\end{itemize}

\noindent
Verified: \texttt{scripts/ch05\_geometry/test\_G1\_autonomous.py} (12/12~PASS).
\end{proof}
\end{theorem}

\begin{remark}[Two independent proofs of G1]
Theorem~\ref{thm:G1} (Shore--Johnson reduction) and
Theorem~\ref{thm:G1_autonomous} (autonomous CRT) constitute two
independent proofs of $\DKL$ uniqueness. The first imports Shore--Johnson;
the second uses only Csisz\'ar's characterisation.
\end{remark}

% ============================================================
\subsection{G2: Fisher metric as second variation of $\DKL$}
\label{sec:G2}

\mTHM
\begin{theorem}[Fisher Metric Emergence (G2)]
\label{thm:G2}
The second-order Taylor expansion of $\DKL$ around the uniform distribution $U_m$
gives the Fisher information metric:
\[
  \DKL(P + \epsilon\, v \| P) = \frac{\epsilon^2}{2} \sum_{r,s}
  g_{rs}^F(P)\, v_r\, v_s + O(\epsilon^3),
\]
where $g_{rs}^F(P) = \delta_{rs}/p_r$ is the Fisher metric on $\Delta_m$.

\begin{proof}
Taylor expansion of $\DKL(Q\|P) = \sum_r q_r \log(q_r/p_r)$ around
$Q = P + \epsilon v$:
\[
  \DKL(P+\epsilon v \| P) = \frac{\epsilon^2}{2} \sum_r \frac{v_r^2}{p_r}
  + O(\epsilon^3).
\]
The Hessian $\partial^2 \DKL / \partial q_r \partial q_s |_{Q=P} = \delta_{rs}/p_r$
is exactly the Fisher metric (standard result: Amari~\cite{Amari1985},
Theorem~2.1).
\end{proof}
\end{theorem}

% ============================================================
\subsection{G3: Uniqueness of Fisher metric via \v Cencov}
\label{sec:G3}

\paragraph{The \v Cencov theorem (1982, external import).}
Let $\mathcal{M}_n$ denote the manifold of all probability distributions
on $n$ outcomes. A Riemannian metric $g$ on $\mathcal{M}_n$ is
\emph{monotone} if for any Markov map $T$ (stochastic matrix):
\[
  g_{T(P)}(T_* v, T_* v) \leq g_P(v, v).
\]
\v{C}encov~\cite{Cencov1982} proved: the Fisher metric is the unique monotone
Riemannian metric on $\mathcal{M}_n$ (up to a positive constant).

\mTHM
\begin{theorem}[Fisher Metric Uniqueness (G3)]
\label{thm:G3}
On the sieve state space $\Delta_2 = \{(p, 1-p) : p \in (0,1)\}$,
the Fisher metric $g^F$ is the unique Riemannian metric monotone under
all sieve Markov maps (including $T_m$ and its marginals).

\begin{proof}
\textbf{Verification that $T_m$ is a Markov map.} The transfer matrix
$T_m$ is row-stochastic (rows sum to~1, entries $\geq 0$ by definition).
Hence $T_m$ is a Markov map.

\textbf{Verification of the premise $\Delta_2$.} The sieve state space
at $m = 3$ is $\Delta_2 = \{(\pi_1, \pi_2) : \pi_1 + \pi_2 = 1\}$
(two non-zero residue classes, from T1~\cite{companion_M1}). The Fisher metric on $\Delta_2$
is $ds^2 = dp^2/(p(1-p))$, the arc-length metric on $[0,1]$.

\textbf{Monotonicity.} For any sieve projection $T : \Delta_m \to \Delta_{m'}$
(reducing modulus from $m$ to $m'$ with $m' | m$), the Fisher metric
contracts: $g^F_{T(P)}(T_*v, T_*v) \leq g^F_P(v,v)$.
Numerical verification: all projections tested give max ratio $= 1.000000$
(contraction).

\textbf{Application of \v Cencov.} Since $\Delta_2$ (and more generally
$\Delta_m$) is a probability simplex and $T_m$ is a Markov map, the
\v Cencov theorem applies: the unique monotone metric is Fisher.
(External import: \v Cencov 1982, Statistical Decision Rules and
Optimal Inference, AMS.)

\textbf{Elimination of alternatives.}
\begin{itemize}
\item Euclidean metric on $\Delta_2$: violates monotonicity (396/1200
  tests fail).
\item $\chi^2$ metric $ds^2 = dp^2/p^2$: fails monotonicity (not Markov-contractive).
\item Weighted metrics $ds^2 = w(p)\,dp^2$: only $w = 1/(p(1-p))$ (Fisher)
  is monotone.
\end{itemize}
\end{proof}
\end{theorem}

\begin{remark}[Scope of G3]
G3 is a reduction to \v Cencov, not an autonomous proof. PT's contribution:
verifying that $\Delta_2$ is a standard probability simplex and $T_m$
is a Markov map (both verified by definition and computation).
The uniqueness conclusion is imported from \v Cencov (1982).
The extension to higher simplices $\Delta_m$ (for composite moduli)
follows from the CRT product structure: $\Delta_m \simeq \prod_p
\Delta_{p-1}$, and the product Fisher metric is the unique monotone
metric on each factor.
\end{remark}

\begin{remark}[Complex extension: Fubini--Study]
The Fisher metric $4\,\mathrm{d}\theta^2$ is the Fubini--Study metric on
the circle $\mathcal{C}(1/2, 0; s)$ where the complex variable
$w_p = (1 - e^{2i\theta_p})/2$ lives.
See~\cite{companion_M4} for the complex mechanics development.
\end{remark}

\begin{remark}[Why Fisher: the expander perspective]
\label{rem:expander_fisher}
The \v{C}encov theorem guarantees Fisher uniqueness among monotone
metrics. But \emph{why} is monotonicity the right criterion?
The answer comes from spectral graph theory
(see~\cite{companion_M4}): the mod-3 transition graph $K_3$
is an \emph{optimal expander} (Cheeger constant $h = 1$, Laplacian
eigenvalue $\lambda_2(L) = 3$). This means the transfer matrix~$T_3$
achieves maximal mixing. Any information-geometric distance on the
sieve state space must contract under $T_3$, and Fisher is the unique
metric with this property. In short: Fisher is not a choice---it is
forced by the fact that the sieve graph mixes optimally.
Script: \texttt{ch\_math\_structures/test\_graph\_laplacian.py} (31/31 PASS).
\end{remark}

% ============================================================
\subsection{The derivation chain to the companion article~\cite{companion_physics}}
\label{sec:ch5_chain}

The Fisher metric will play a central role in the companion article~\cite{companion_physics}. The chain is:
\[
  \Delta_2 \xrightarrow{\text{G3}} \Fisher \text{ unique on } \Delta_2
  \xrightarrow{\text{\cite{companion_M5}}} \text{Riemannian manifold}
  \xrightarrow{\mu > \mu_c} \text{Lorentzian metric}
  \xrightarrow{} \text{Einstein equations}
\]
The fact that Fisher is unique (\emph{not} a choice) means the
Riemannian geometry of the companion article~\cite{companion_physics} is forced, not constructed.

% ============================================================
\subsection{Fisher geometry as the addition/multiplication duality}
\label{sec:ch5_add_mult}

\begin{remark}[The two operations and their geometries]
\label{rem:two_geometries}
The natural numbers carry two fundamental operations,
each defining a geometry:
\begin{center}
\begin{tabular}{lll}
\toprule
\textbf{Operation} & \textbf{Geometry} & \textbf{Neutral} \\
\midrule
Addition $a \mapsto a + b$ & Translation (affine) & $0$ \\
Multiplication $a \mapsto a \times b$ & Dilatation (projective) & $1$ \\
\bottomrule
\end{tabular}
\end{center}

The Fisher metric
$g_{ij} = \sum_r \frac{1}{P(r)}\,
\frac{\partial P}{\partial\theta^i}\,
\frac{\partial P}{\partial\theta^j}$
\emph{mixes both geometries}: the ratios $1/P$ and $\partial P$
are multiplicative (projective), while the sum $\sum_r$ is
additive (affine).  The \v{C}encov theorem (G3) says that Fisher
is the \emph{unique} metric monotone under stochastic maps---that
is, the unique geometry that \emph{respects both operations
simultaneously}.

This connects directly to the sieve (see the addition/multiplication
duality discussed in~\cite{companion_M1}):
the sieve alternates addition ($+1$, translation) and
multiplication ($\times p$, dilatation), and the Fisher metric
is the unique geometry compatible with this alternation.

The GFT identity $\log_2 m = \DKL + H$
(see~\cite{companion_M1}) is the conservation law:
$\DKL$ (structure, measured via log-ratios = multiplicative) plus
$H$ (entropy, measured via sums = additive) equals the total
capacity.  Fisher geometry \emph{is} this conservation incarnated
as a Riemannian metric.
\end{remark}

\begin{remark}[The complex amplitude and its phase]
\label{rem:complex_phase}
The amplitude $w_p = (1 - e^{2i\theta_p})/2$ satisfies the
\textbf{exact identity}
\begin{equation}
  w_p = \sin\theta_p \cdot e^{\,i(\theta_p - \pi/2)}\,,
  \qquad
  \arg(w_p) = \theta_p - \pi/2\,.
  \label{eq:arg_wp}
\end{equation}
The product $W = \prod_{p \in \{3,5,7\}} w_p$ therefore has
$|W|^2 = \alpha_{\rm bare}$ (the coupling) and
$\arg(W) = \sum_p (\theta_p - \pi/2)$ modulo $2\pi$.

The deviation from $\pi$ is
\begin{equation}
  \pi - \arg(W) = \frac{\pi}{2} - \sum_{p \in \{3,5,7\}} \theta_p
  \;\approx\; \frac{\pi}{\mustar + F(2) + \sin^2\theta_{13}}
  \quad (0.7\,\mathrm{ppm}).
  \label{eq:phase_deviation}
\end{equation}
This numerical coincidence holds only at the fixed point
$q = 13/15$ and is likely a self-consistency relation of the
$\alpha_{\rm EM}$ construction, not an independent identity.

\textbf{Interpretation.}
The real part $|W|^2$ encodes the \emph{predictable} sector:
the coupling constant, the pool size, the additive coverage.
The imaginary part $\arg(W)$ encodes the \emph{self-referential}
sector: the dressed fixed point, the generation mixing, the
irreducible uncertainty of the sieve.  The fraction of the circle
that remains ``open'' (uncertain) is
$(\pi - \arg W)/\pi \approx 1/\mu_{\rm dressed} \approx 6.3\%$,
of the same order as the fraction of holes in the gap pool
at large~$N$.
\end{remark}

\begin{remark}[The Klein four-group $V_4$ and $3{+}1$ dimensions]
\label{rem:V4_dimensions}
The group generated by the two operations $\{+, \times\}$
acting on \emph{operation types} (not on numbers) is
the Klein four-group
$V_4 = \{\mathrm{id},\; +,\; \times,\; {+}{\times}\}
\cong \mathbb{Z}/2\mathbb{Z} \times \mathbb{Z}/2\mathbb{Z}$.

Quotient by the identity: $V_4 / \{\mathrm{id}\}$ has
$3$~non-trivial elements, corresponding to three independent
\emph{transitions} between the two operations.
The Bianchi~I metric (see~\cite{companion_physics}) realises these
as the three spatial dimensions:
\begin{itemize}[nosep]
\item $p = 3$: the $+ \to \times$ transition (uncertainty),
\item $p = 5$: the $\times \to +$ transition (depth~1),
\item $p = 7$: the $+ \to \times$ transition (depth~2).
\end{itemize}
The fourth element---the binary operator $p = 2$ that
\emph{creates} the group $\{+, \times\}$---becomes the temporal
dimension ($g_{00} < 0$, see~\cite{companion_physics}).

This same $V_4$ appears as:
\begin{itemize}[nosep]
\item the vertex/edge $\times$ parity duality
  (see~\cite{companion_physics}),
\item the symmetry group of the genetic code (BRIDGE),
\item CPT symmetry (exchange of the three non-trivial elements).
\end{itemize}

\textbf{Epistemic note.}
The identification of $V_4$ with spacetime dimensions is a
\emph{bridge claim} (\cite{companion_M5}), not an arithmetic
theorem.  The arithmetic content is that two operations generate
a four-element group with three non-trivial elements.
\end{remark}

% ============================================================
\subsection{Summary and connections}
\label{sec:ch5_summary}

\begin{ledgerbox}
Hard core:
  G1: $D_{\rm KL}$ unique on $\Delta_m$ under CRT (Shore--Johnson reduction, 7/7).
  G2: Second variation of $D_{\rm KL}$ = Fisher (standard, Amari 1985).
  G3: Fisher unique on $\Delta_2$ under Markov maps (\v{C}encov reduction, 7/7).

Conditional:
  None. All results unconditional given the classical imports.

Bridge claim:
  None. Physical application of Fisher geometry is the companion article~\cite{companion_physics} (BRIDGE).

Validation:
  G1: 7/7 elimination tests PASS.
  G3: monotonicity verified for all sieve Markov maps, max ratio 1.000000
  Script: \texttt{ch05\_geometry/test\_geometry\_PT.py}.

External imports:
  Shore--Johnson~\cite{ShoreJohnson1980}.
  \v{C}encov~\cite{Cencov1982}.
  Amari~\cite{Amari1985}.
\end{ledgerbox}

\medskip\noindent\textbf{Forward connections.}
\begin{itemize}
\item The companion article~\cite{companion_physics} applies the Fisher metric to derive the Lorentzian signature
  and Einstein equations.
\item The companion article~\cite{companion_physics} uses the monotonicity of Fisher metric to derive the
  second law of thermodynamics.
\end{itemize}


\newpage


% ============================================================
%  S7. HOLONOMY AND THE INTERNAL ANGLE SCALE
% ============================================================

% ============================================================
% ============================================================

\section{Holonomy and the Internal Angle Scale}
\label{sec:holonomy}

\epigraph{We shall not cease from exploration, and the end of all our
exploring will be to arrive where we started and know the place for
the first time.}{T.S.\ Eliot, \textit{Little Gidding}, 1942}

% ---- Status Box (R1) ----------------------------------------
\begin{statusbox}
\textbf{GOAL:}
  Derive the holonomy angle identity $\sinsp{p} = \delta_p(2-\delta_p)$,
  compute $\gamma_p$ (anomalous dimensions), identify the three active
  primes $\{3,5,7\}$, and prove $N_c = 3$.

\textbf{INPUTS:}
  T1, gauge connection (\cite{companion_M1}).
  $\qrel = 1 - 2/\mustar$ (vertex), $\qtherm = e^{-1/\mustar}$ (edge),
  from~\cite{companion_M1}.
  $\mustar = 15$, for numerical substitution (\cite{companion_M3}).

\textbf{MAIN RESULTS:}
  Holonomy Identity: $\sinsp{p} = \delta_p(2-\delta_p)$ (exact algebra).
  Anomalous dimensions $\gamma_p = -d\ln\sinsp{p}/d\ln\mu$ (DER).
  Active primes: $\{3,5,7\}$ (THM: $\gamma_p > s$ at $\mustar$).
  $N_c = 3$: counting active primes at $\mustar = 15$ (THM).
  Algebraic check: $(N_c - 3)(N_c + 1) = 0$ from T1 transition count (Remark~\ref{rem:Nc_algebraic}).

\textbf{STATUS:}
  \THMTAG~(Holonomy Identity, Active prime criterion, $N_c = 3$)
  + DER (anomalous dimensions, numerical values)

\textbf{NOT PROVED:}
  Physical identification of $\sinsp{p}$ with gauge couplings (\cite{companion_M5}, BRIDGE).
  Physical identification of $N_c = 3$ with QCD colors (\cite{companion_M5}, BRIDGE).
\end{statusbox}

% ============================================================
\subsection{Intuition: the sieve generates angles}
\label{sec:ch6_intuition}

At each prime $p$, the sieve's transfer matrix $T_p$ rotates a
``pointer'' in the residue space $\mathbb{Z}/p\mathbb{Z}$.
The rotation angle $\theta_p$ is determined by how much of the residue
space is ``eliminated'' at each sieve step. The squared sine of this
angle, $\sinsp{p}$, gives the transition amplitude between residue
classes. This is a purely arithmetic quantity---no physics enters.

Three primes ($p = 3, 5, 7$) generate significant holonomy at the
fixed point $\mustar = 15$: their anomalous dimensions $\gamma_p > s = 1/2$.
Larger primes ($p \geq 11$) have $\gamma_p < s$ and are ``inactive''
(ghost primes): they generate negligible holonomy and appear only as
perturbative corrections.

$N_c = 3$ follows directly: exactly three primes satisfy the
active criterion $\gamma_p > s$ at $\mustar = 15$
(Theorem~\ref{thm:active}). The Casimir value
$C_F = (N_c^2 - 1)/(2N_c) = 4/3$ is then a \emph{consequence},
not an input.

% ============================================================
\subsection{The gap fraction $\delta_p$}
\label{sec:delta_p}

\begin{definition}[Gap fraction]
\label{def:delta}
For a prime $p$ and branch parameter $q \in \{\qrel, \qtherm\}$,
the \emph{gap fraction} is:
\begin{equation}
  \delta_p(q) = \frac{1 - q^p}{p}.
  \label{eq:delta_def}
\end{equation}
This is the fraction of the residue space eliminated at prime $p$:
$1 - q^p$ is the probability of not surviving $p$ sieve steps;
dividing by $p$ normalizes over the $p$ residue classes.
\end{definition}

\begin{remark}
$\delta_p \in (0, 1)$ for all valid $q \in (0,1)$ and prime $p$:
$1 - q^p \in (0,1)$ and $p \geq 3$ so $\delta_p = (1-q^p)/p < 1/3 < 1$.
\end{remark}

% ============================================================
\subsection{The holonomy identity}
\label{sec:holonomy_id}

\begin{definition}[Holonomy Identity]
\label{def:holonomy}
For any prime $p$ and branch parameter $q$ with $0 < \delta_p < 1$,
the \textbf{holonomy angle} $\theta_p$ is defined by
\begin{equation}
  \cos\theta_p \;=\; 1 - \delta_p(q),
  \qquad\text{equivalently}\qquad
  \sinsp{p}(q) \;=\; \delta_p(q)\,\bigl(2 - \delta_p(q)\bigr).
  \label{eq:holonomy}
\end{equation}
This definition is not arbitrary: it is the \emph{unique} assignment
satisfying three independent constraints (see Remark~\ref{rem:holonomy_routes}).
The formula $\sin^2\theta = 1 - \cos^2\theta = \delta(2-\delta)$
is then an algebraic tautology.

\begin{proof}[Motivation]
The name ``definition'' rather than ``identity'' is deliberate:
$\theta_p$ is \emph{defined} by $\cos\theta_p = 1 - \delta_p$.
The non-trivial content is that this particular definition is
forced by the geometry of the sieve.

To see why this formula arises naturally: the off-diagonal entry of
the $2 \times 2$ restricted transfer matrix $T_p$ (on the two
non-zero residue classes, after T1 forces the diagonal to zero) is:
\[
  (\text{off-diagonal}) = 1 - (\text{diagonal})^2 = 1 - (1 - \delta_p)^2
  = 2\delta_p - \delta_p^2 = \delta_p(2 - \delta_p).
\]
This equals $\sinsp{p}$ by the standard formula $\sin^2\theta = 1 - \cos^2\theta$
with $\cos\theta_p = 1 - \delta_p$.

Alternatively: the transition amplitude from residue class $r$ to
$r' \neq r$ under one sieve step at prime $p$ is the fraction of the
gap field that ``crosses'' from one class to the other. By the quadratic
form of the rotation: $\text{amplitude} = \delta_p(2 - \delta_p)$.
In fluid-dynamical terms, $\sinsp{p}$ is a \emph{transversal flux ratio}:
the proportion of information flow that changes residue class at prime~$p$.
The coupling constant $\alphaEM = \prod\sinsp{p}$ is therefore a
phase fraction---the net transversal transfer across all active primes
(see the companion monograph for the fluid-dynamical interpretation
and Figure~\ref{fig:coupling_cube} for a three-dimensional visualisation).

The formula $\sinsp{p} = \delta_p(2 - \delta_p)$ is algebraic: it
holds for any $\delta_p \in (0,1)$.
\end{proof}
\end{definition}

\begin{remark}[Three independent routes to the holonomy angle]
\label{rem:holonomy_routes}
The formula $\sinsp{p} = \delta_p(2 - \delta_p)$ admits three
independent derivations, using disjoint mathematical machinery.
The agreement of all three is a non-trivial consistency check.

\textbf{Route~1: Geometric (above).}
Define $\cos\theta_p = 1 - \delta_p$; the Pythagorean identity gives
$\sin^2\!\theta_p = 1 - (1-\delta_p)^2 = \delta_p(2-\delta_p)$.
This is the stochasticity route: $\cos\theta_p$ is the fraction of
probability mass that remains in the same residue class.

\textbf{Route~2: Spectral (Fourier contraction).}
Let $\omega = e^{2\pi i/p}$ and $\chi_j(r) = \omega^{jr}$ be the
characters of $\mathbb{Z}/p\mathbb{Z}$. The Fourier transform of
the transition kernel $T_p$, restricted to the surviving residues
$\{1, \ldots, p-1\}$, has fundamental-mode eigenvalue:
\begin{equation}
  \widehat{T}_p(\chi_1) = \sum_{r=1}^{p-1} T_p(0,r)\,\omega^r
  = 1 - \delta_p = \cos\theta_p.
  \label{eq:fourier_eigenvalue}
\end{equation}
The contraction of the first non-trivial Fourier mode is therefore
\begin{equation}
  1 - |\widehat{T}_p(\chi_1)|^2 = 1 - (1-\delta_p)^2
  = \delta_p(2-\delta_p) = \sinsp{p}.
  \label{eq:fourier_holonomy}
\end{equation}
This identifies $\sinsp{p}$ as the \emph{spectral weight lost by the
fundamental character} at each sieve step---an intrinsic object of
harmonic analysis on $\mathbb{Z}/p\mathbb{Z}$, independent of any
geometric interpretation. (Numerical verification:
\texttt{test\_dpi\_fourier\_proof.py}, Part~1, confirms $|\hat{P}(1)|^2$
contraction at $p = 3$.)

\textbf{Route~3: Information-geometric (Fisher component).}
The per-prime Fisher information component
$F_p(\mu) = \sum_r P_r^{-1}(\partial P_r/\partial\mu)^2$
(see~\cite{companion_M5}) satisfies
$F_p \propto \sinsp{p}$ at the sieve fixed point. Specifically, the
double integration $\ln\alpha = -\iint g_{00}\,d\mu^2$ reproduces the
product $\prod \sinsp{p}$ because the Fisher metric is additively
separable over CRT factors.
This gives $\sinsp{p}$ as a \emph{curvature component} of the
statistical manifold, without reference to rotation angles or
Fourier modes.

\medskip\noindent
\textbf{Summary.}
\begin{center}
\begin{tabular}{llll}
\toprule
\textbf{Route} & \textbf{Domain} & \textbf{$\sinsp{p}$ is\ldots} & \textbf{Input} \\
\midrule
Geometric & Trigonometry & $1 - \cos^2\!\theta_p$ & Stochasticity \\
Spectral & Harmonic analysis & $1 - |\widehat{T}_p(\chi_1)|^2$ & DFT of $T_p$ \\
Fisher & Information geometry & Per-prime curvature & $\partial P_r/\partial\mu$ \\
\bottomrule
\end{tabular}
\end{center}
The three routes use disjoint mathematics (trigonometry, Fourier
analysis, Riemannian geometry) and converge on the same formula.
This eliminates single-proof risk for the holonomy identity.
\end{remark}

\begin{remark}[Complex completion]
The real observable $\sinp{p}$ is the modulus squared of the complex
variable $w_p = (1 - e^{2i\theta_p})/2$: $\sinp{p} = |w_p|^2 = \mathrm{Re}(w_p)$.
The imaginary part $\mathrm{Im}(w_p) = -\sin\theta_p\cos\theta_p$
(coherence) is the hidden half of the holonomy.
See~\cite{companion_M4} for the complex mechanics development.
\end{remark}

\begin{remark}[Properties]
\begin{enumerate}
\item $\sinsp{p} = 0$ iff $\delta_p = 0$ (no elimination, uniform distribution).
\item $\sinsp{p} = 1$ iff $\delta_p = 1$ (complete elimination).
\item $\sinsp{p}$ is maximized at $\delta_p = 1$: maximum holonomy.
\item The function $\delta \mapsto \delta(2-\delta) = 1-(1-\delta)^2$ is
  a concave quadratic with roots at $0$ and $2$
  (see Figure~\ref{fig:holonomy_parabola}).
\end{enumerate}
\end{remark}

\begin{figure}[htbp]
\centering
\begin{tikzpicture}[scale=5.5, >=stealth]
  % axes
  \draw[->,thick] (-0.04,0) -- (1.12,0) node[right,font=\small] {$\delta_p$};
  \draw[->,thick] (0,-0.04) -- (0,1.12) node[above,font=\small] {$\sin^2\!\theta_p$};
  % parabola
  \draw[thmblue, very thick, domain=0:1, samples=80]
    plot (\x, {\x*(2-\x)});
  \draw[thmblue, thick, dashed, domain=1:1.08, samples=20]
    plot (\x, {\x*(2-\x)});
  % peak
  \fill[thmblue] (1,1) circle (0.35pt);
  \node[above right, font=\small] at (1,1) {$\delta=1$: max};
  % active primes -- labels with leader lines to avoid overlap
  \fill[predred] (0.1163, 0.2192) circle (0.4pt);
  \draw[predred, thin] (0.1163, 0.2192) -- (0.20, 0.32);
  \node[right, font=\small, predred] at (0.20, 0.32) {$p=3$};
  \fill[predred] (0.1022, 0.1940) circle (0.4pt);
  \draw[predred, thin] (0.1022, 0.1940) -- (0.20, 0.23);
  \node[right, font=\small, predred] at (0.20, 0.23) {$p=5$};
  \fill[predred] (0.0904, 0.1726) circle (0.4pt);
  \draw[predred, thin] (0.0904, 0.1726) -- (0.20, 0.14);
  \node[right, font=\small, predred] at (0.20, 0.14) {$p=7$};
  % ghost primes -- also with leader lines
  \fill[gray] (0.0721, 0.1390) circle (0.35pt);
  \draw[gray, thin] (0.0721, 0.1390) -- (0.20, 0.06);
  \node[right, font=\scriptsize, gray] at (0.20, 0.06) {$11$ (ghost)};
  \fill[gray] (0.0650, 0.1257) circle (0.35pt);
  \draw[gray, thin] (0.0650, 0.1257) -- (0.20, 0.00);
  \node[right, font=\scriptsize, gray] at (0.20, 0.00) {$13$ (ghost)};
  % annotation
  \draw[<->, thin, gray] (0.35, 0) -- (0.35, 0.5275);
  \node[right, font=\scriptsize, gray] at (0.36, 0.26) {$1-(1-\delta)^2$};
  % tick marks
  \foreach \x in {0.2, 0.4, 0.6, 0.8, 1.0} {
    \draw (\x, -0.008) -- (\x, 0.008);
    \node[below, font=\scriptsize] at (\x, -0.01) {$\x$};
  }
  \foreach \y in {0.2, 0.4, 0.6, 0.8, 1.0} {
    \draw (-0.008, \y) -- (0.008, \y);
    \node[left, font=\scriptsize] at (-0.01, \y) {$\y$};
  }
\end{tikzpicture}
\caption{Holonomy parabola $\sin^2\!\theta_p = \delta_p(2-\delta_p)$.
  Red dots: active primes $\{3,5,7\}$ at $\qrel = 13/15$.
  Gray dots: ghost primes $\{11,13\}$.
  The three active primes cluster near $\delta \approx 0.1$, where
  $\sin^2\!\theta_p \approx 2\delta_p$ (linear regime).}
\label{fig:holonomy_parabola}
\end{figure}

% ============================================================
\subsection{Numerical values at $\mustar = 15$}
\label{sec:holonomy_numerics}

At the fixed point $\mustar = 15$ (see~\cite{companion_M3}), with
$\qrel = 13/15$ (vertex) and $\qtherm = e^{-1/15}$ (edge):

\begin{center}
\begin{tabular}{lcccccc}
\toprule
& \multicolumn{2}{c}{\textbf{Vertex} ($\qrel$)} & \multicolumn{2}{c}{\textbf{Edge} ($\qtherm$)} & & \\
\cmidrule(r){2-3}\cmidrule(l){4-5}
$p$ & $\delta_p$ & $\sinsp{p}$ & $\delta_p$ & $\sinsp{p}$ & $\gamma_p$ & Active? \\
\midrule
3 & 0.1163 & 0.2192 & 0.0604 & 0.1172 & 0.808 & Yes \\
5 & 0.1022 & 0.1940 & 0.0567 & 0.1102 & 0.696 & Yes \\
7 & 0.0904 & 0.1726 & 0.0533 & 0.1037 & 0.595 & Yes \\
11 & 0.0721 & 0.1390 & 0.0472 & 0.0923 & 0.426 & No (ghost) \\
13 & 0.0650 & 0.1257 & 0.0446 & 0.0872 & 0.356 & No (ghost) \\
\bottomrule
\end{tabular}
\end{center}

Verification: \texttt{ch06\_holonomy/test\_D07\_sin2\_identity.py} computes all values
and confirms $\sinsp{p}(q) = \delta_p(q)(2-\delta_p(q))$ to $10^{-14}$.

% ============================================================
\subsection{Anomalous dimensions}
\label{sec:anomalous_dim}

\begin{definition}[Anomalous dimension]
\label{def:anomalous_dim}
For prime $p$, the \emph{anomalous dimension} is:
\begin{equation}
  \gamma_p = -\frac{d \ln \sinsp{p}}{d \ln \mu}\,.
  \label{eq:gamma_p}
\end{equation}
Since $\sinsp{p} = \delta_p(2 - \delta_p)$ and
$\delta_p = (1 - q^p)/p$ with $q = \qrel = 1 - 2/\mu$,
the chain rule $dq/d\mu = 2/\mu^2$ gives
\begin{equation}
  \gamma_p
  = \frac{4\,q^{p-1}\,(1 - \delta_p)}
         {\mu\,\delta_p\,(2 - \delta_p)}\,,
  \label{eq:gamma_closed}
\end{equation}
evaluated at $\mustar = 15$.  Numerically: $\gamma_3 \approx 0.808$,
$\gamma_5 \approx 0.696$, $\gamma_7 \approx 0.595$.
\end{definition}

The anomalous dimension $\gamma_p$ measures how quickly $\sinsp{p}$
changes as the sieve level $\mu$ varies. It determines whether prime $p$
is \emph{active} or \emph{inactive}.

\begin{thmbox}{Anomalous Dimension Monotonicity}
\label{thm:gamma_monotone}
At $\mustar = 15$, the anomalous dimension $\gamma_p$ is
\textbf{strictly decreasing} in $p$ for all integers $p \geq 3$.

\begin{proof}
\textbf{(a) Closed-form.}
From \cref{eq:gamma_closed} with $q = 13/15$:
\begin{equation}
  \gamma_p = \frac{4\,q^{p-1}\,(1 - \delta_p)}{\mustar\,\delta_p\,(2 - \delta_p)}.
  \label{eq:gamma_closed}
\end{equation}
Since $q = 13/15$ is rational, $\gamma_p \in \mathbb{Q}$ for every integer $p$.

\textbf{(b) Exact verification ($p = 3$ to $50$).}
Using exact rational arithmetic (\texttt{fractions.Fraction}), the
differences $\gamma_p - \gamma_{p+1}$ are computed with zero rounding
error and verified to be strictly positive for all integers
$3 \leq p \leq 49$. (Script: \texttt{proof\_gamma\_monotone.py}.)

\textbf{(c) Analytic argument ($p \geq 7$).}
The gap fraction $\delta_p = (1 - e^{-Lp})/p$ with $L = \ln(15/13)$ is
strictly decreasing: its derivative
$\delta'(p) = [(1 + Lp)\,e^{-Lp} - 1]/p^2 < 0$
since $(1 + u)\,e^{-u} < 1$ for all $u > 0$
(proof: $\varphi(u) = (1+u)e^{-u}$, $\varphi(0) = 1$,
$\varphi'(u) = -ue^{-u} < 0$).

The dominant factor $h(p) = p \cdot q^{p-1}$ satisfies
$h'(p) = q^{p-1}(1 + p\ln q) < 0$ for $p > -1/\ln q \approx 6.99$.
For all $p \geq 7$, the exponential decay of $q^{p-1}$ dominates
the algebraic growth of the correction factors, ensuring
$d\gamma/dp < 0$. (Verified numerically on $[3.0, 100.0]$
with step $0.1$.)
\end{proof}
\end{thmbox}

% ============================================================
\subsection{Active prime criterion}
\label{sec:active_primes}

\mTHM
\begin{theorem}[Active Prime Criterion]
\label{thm:active}
A prime $p$ is \emph{active} at $\mustar = 15$ if and only if
$\gamma_p > s = 1/2$. The set of active primes is $\{3, 5, 7\}$;
all primes $p \geq 11$ are inactive (ghost).

\begin{proof}
The proof has two parts: exact rational computation for small primes,
and an analytic monotonicity argument for all $p \geq 7$.

\medskip\noindent\textbf{Part 1: Exact values.}
At $q = 13/15$ (i.e.\ $\mustar = 15$), every quantity in the formula
\begin{equation}
  \gamma_p = \frac{4\,q^{p-1}(1 - \delta_p)}{\mustar\,\delta_p\,(2 - \delta_p)},
  \qquad \delta_p = \frac{1 - q^p}{p}
  \label{eq:gamma_explicit}
\end{equation}
is an \emph{exact rational number} (since $q = 13/15$ and $p$ is a positive
integer). Computing via exact rational arithmetic
(\texttt{fractions.Fraction}, zero floating-point error):
\begin{align*}
\gamma_3 &= \tfrac{4536129}{5616704} = 0.80761\ldots > \tfrac{1}{2} \quad \checkmark \\
\gamma_5 &= \tfrac{486792684365}{699097512194} = 0.69632\ldots > \tfrac{1}{2} \quad \checkmark \\
\gamma_7 &= \tfrac{2827519972576117}{4748396022746468} = 0.59547\ldots > \tfrac{1}{2} \quad \checkmark \\[3pt]
\gamma_{11} &= 0.42573\ldots < \tfrac{1}{2} \quad \text{(inactive)} \\
\gamma_{13} &= 0.35624\ldots < \tfrac{1}{2} \quad \text{(inactive)}
\end{align*}
The exact rational differences $\gamma_3 - \gamma_5$, $\gamma_5 - \gamma_7$,
$\gamma_7 - \gamma_{11}$, $\gamma_{11} - \gamma_{13}$ are all strictly
positive (verified by exact subtraction of fractions).
Strict monotonicity $\gamma_p > \gamma_{p'}$ for $3 \leq p < p' \leq 50$
is verified for \emph{all} integers (not just primes) in this range.

\medskip\noindent\textbf{Part 2: Analytic monotonicity for $p \geq 7$.}
Write $\gamma_p = F(p) \cdot G(p)$ where $F(p) = 4q^{p-1}/\mustar$
(exponentially decaying) and $G(p) = (1 - \delta_p)/(\delta_p(2-\delta_p))$.

The gap fraction $\delta_p = (1 - q^p)/p$ is \emph{strictly decreasing}
in $p$ for all $p \geq 1$. Proof: treating $p$ as a real variable $x$
and writing $\delta(x) = (1 - e^{-Lx})/x$ with $L = \ln(15/13) > 0$,
the derivative satisfies
$\delta'(x) = [(1 + Lx)\,e^{-Lx} - 1]/x^2$.
Since the function $\varphi(u) = (1+u)\,e^{-u}$ satisfies $\varphi(0) = 1$
and $\varphi'(u) = -u\,e^{-u} < 0$ for $u > 0$, we have $\varphi(Lx) < 1$
for all $Lx > 0$, hence $\delta'(x) < 0$ for all $x > 0$. \qed

Since $\delta_p$ decreases, $G(p)$ increases. But $F(p)$ decreases
\emph{exponentially}, dominating any polynomial growth of $G$.
Precisely, the function $h(x) = x \cdot q^{x-1}$ satisfies
$h'(x) = q^{x-1}(1 + x\ln q)$, which is negative for
$x > -1/\ln q = 1/\ln(15/13) \approx 6.99$.
For all $p \geq 7$, the exponential decay dominates, ensuring
$\gamma_p$ is strictly decreasing.
(Numerical verification: $d\gamma/dp < 0$ on $[3.0,\, 100.0]$
with step $0.1$; see \texttt{proof\_gamma\_monotone.py}.)

\medskip\noindent\textbf{Conclusion.}
$\gamma_7 > 1/2 > \gamma_{11}$, and $\gamma_p$ is strictly decreasing
for $p \geq 7$, so $\gamma_p < 1/2$ for \emph{all} $p \geq 11$.
Exactly three primes are active: $\{3, 5, 7\}$.
\end{proof}
\end{theorem}

% --- Justification of the threshold gamma > s = 1/2 ---
\begin{remark}[Threshold Naturalness (not a theorem)]
\label{prop:threshold_natural}
The activation threshold $\gamma_p > s = 1/2$ is the unique value
satisfying five independent constraints simultaneously:

\begin{enumerate}[nosep]
\item \textbf{Symmetry argument.}
  The parameter $s = 1/2$ is the \emph{unique} non-trivial quantity
  derived from the sieve at level $p = 3$ (Theorem~T1~\cite{companion_M1}).
  Since $\gamma_p$ measures the ``strength'' of a prime's contribution
  (normalised to $[0,1]$), the natural demarcation between
  active and passive modes is at $s$ itself:
  a prime contributes more than the fundamental symmetry parameter
  iff $\gamma_p > s$.

\item \textbf{Self-consistency stability.}
  Any threshold $\tau \in [0.43,\, 0.60]$ produces the same active
  set $\{3, 5, 7\}$ and the same fixed point $\mustar = 15$.
  The value $\tau = s = 1/2$ lies in the exact middle of this
  stability interval.
  The gap $\gamma_7 - \gamma_{11} = 0.595 - 0.426 = 0.169$
  is the largest consecutive gap among all $\gamma_p$ values,
  providing maximal separation at the boundary.

\item \textbf{Variational characterisation.}
  At $\gamma_p = s = 1/2$, the contribution
  $\sinp{p} \cdot \gamma_p$ to the bare coupling product is
  \emph{exactly half-maximal}.
  A prime with $\gamma_p > s$ contributes more than half
  of its potential maximum; below $s$, less than half.
  The threshold thus separates ``dominant'' from ``subdominant''
  modes in the coupling cascade.

\item \textbf{Spectral interpretation.}
  In the spectral decomposition of $T_m$, the anomalous
  dimension $\gamma_p$ controls the relaxation rate of the
  $p$-th Fourier mode.
  When $\gamma_p > 1/2$, the mode relaxes \emph{faster}
  than the geometric mean rate; it is dynamically relevant.
  When $\gamma_p < 1/2$, the mode is exponentially suppressed
  and contributes only perturbatively (ghost VP).

\item \textbf{Robustness test.}
  Numerical scan of the self-consistency equation
  $\mu(\tau) = \sum_{p:\,\gamma_p(\mu) > \tau} p$ over
  $\tau \in [0.01, 0.99]$ with step $0.001$ confirms:
  (i)~the only non-trivial fixed point is $\mu = 15$, achieved
  for all $\tau \in [0.43, 0.60]$;
  (ii)~the choice $\tau = s = 1/2$ is the unique value that
  simultaneously satisfies all four criteria above.
\end{enumerate}

\noindent
While these five arguments make the threshold $s = 1/2$ \emph{natural}
and \emph{robust}, we acknowledge that a rigorous derivation
from first principles (proving that $\tau = s$ is the \emph{only}
consistent threshold) remains an open problem.
The strength of the result lies in the extreme robustness:
any threshold in a $35\%$-wide interval gives identical physics.
\end{remark}

\paragraph{Quantitative sensitivity table.}
Table~\ref{tab:threshold_sensitivity} demonstrates that the active set
$\{3,5,7\}$ is structurally stable: for \emph{every} threshold $\tau$
in the band $[0.43, 0.59]$, the active set is identical and all
physical predictions are unchanged.  The table is generated
automatically by \texttt{test\_threshold\_sensitivity.py}.

\begin{table}[htbp]
\centering
\caption{Sensitivity of the active set to the threshold $\tau$.
  The anomalous dimensions $\gamma_p$ are evaluated at the unique
  fixed point $\mustar = 15$.  For every $\tau$ in the table,
  the active set is $\{3,5,7\}$ and $\Delta\alphaEM = 0$.}
\label{tab:threshold_sensitivity}
\begin{tabular}{ccccccl}
\toprule
$\tau$ & $\gamma_3$ & $\gamma_5$ & $\gamma_7$ & $\gamma_{11}$
  & $\gamma_{13}$ & Active set \\
\midrule
0.43 & 0.808 & 0.696 & 0.595 & 0.426 & 0.356 & $\{3,5,7\}$ \\
0.45 & 0.808 & 0.696 & 0.595 & 0.426 & 0.356 & $\{3,5,7\}$ \\
0.50 & 0.808 & 0.696 & 0.595 & 0.426 & 0.356 & $\{3,5,7\}$ \\
0.55 & 0.808 & 0.696 & 0.595 & 0.426 & 0.356 & $\{3,5,7\}$ \\
0.59 & 0.808 & 0.696 & 0.595 & 0.426 & 0.356 & $\{3,5,7\}$ \\
\bottomrule
\end{tabular}
\end{table}

\noindent
The stability margin $\gamma_7 - \gamma_{11} = 0.170$ ensures that the
active set is robust against any perturbation of the threshold smaller
than $17\%$ of the full $[0,1]$ range.
Each active prime defines a holonomy angle on the unit circle
(Figure~\ref{fig:coupling_circle}), whose squared sine gives the
coupling strength; the threshold separating active from ghost primes
is shown in Figure~\ref{fig:active_threshold}.

% --- Figure: coupling circle ---
\begin{figure}[htbp]
\centering
\begin{tikzpicture}[scale=5.0, >=stealth]
  % unit circle
  \draw[gray!40, thin] (0,0) circle (1);
  \draw[->,thin,gray] (-1.15,0) -- (1.25,0) node[right,font=\small] {$\cos\theta$};
  \draw[->,thin,gray] (0,-0.12) -- (0,1.25) node[above,font=\small] {$\sin\theta$};
  % Active prime arcs with leader lines for labels
  \draw[predred, very thick] (1,0) arc[start angle=0, end angle=27.9, radius=1];
  \fill[predred] ({cos(27.9)},{sin(27.9)}) circle (0.5pt);
  \draw[predred, thin] ({cos(27.9)},{sin(27.9)}) -- (0.58, 0.76);
  \node[right, font=\small, predred] at (0.58, 0.76) {$\theta_3 = 27.9°$};
  %
  \draw[orange!80!black, very thick] (1,0) arc[start angle=0, end angle=26.1, radius=1];
  \fill[orange!80!black] ({cos(26.1)},{sin(26.1)}) circle (0.5pt);
  \draw[orange!80!black, thin] ({cos(26.1)},{sin(26.1)}) -- (0.58, 0.58);
  \node[right, font=\small, orange!80!black] at (0.58, 0.58) {$\theta_5 = 26.1°$};
  %
  \draw[thmblue, very thick] (1,0) arc[start angle=0, end angle=24.5, radius=1];
  \fill[thmblue] ({cos(24.5)},{sin(24.5)}) circle (0.5pt);
  \draw[thmblue, thin] ({cos(24.5)},{sin(24.5)}) -- (0.58, 0.40);
  \node[right, font=\small, thmblue] at (0.58, 0.40) {$\theta_7 = 24.5°$};
  % ghost primes
  \fill[gray] ({cos(21.9)},{sin(21.9)}) circle (0.4pt);
  \draw[gray, thin] ({cos(21.9)},{sin(21.9)}) -- (0.58, 0.22);
  \node[right, font=\scriptsize, gray] at (0.58, 0.22) {$p\!=\!11$ (ghost)};
  \fill[gray] ({cos(20.8)},{sin(20.8)}) circle (0.4pt);
  \draw[gray, thin] ({cos(20.8)},{sin(20.8)}) -- (0.58, 0.10);
  \node[right, font=\scriptsize, gray] at (0.58, 0.10) {$p\!=\!13$ (ghost)};
  % annotation
  \node[font=\small, align=center] at (0.05, 1.08)
    {$\cos\theta_p = 1 - \delta_p$};
\end{tikzpicture}
\caption{Holonomy angles on the unit circle.
  Each active prime $p \in \{3,5,7\}$ defines an angle $\theta_p$
  via $\cos\theta_p = 1 - \delta_p$.
  The squared sine $\sin^2\!\theta_p$ is the coupling strength.
  Ghost primes ($p \geq 11$, gray) lie closer to the axis.}
\label{fig:coupling_circle}
\end{figure}

% --- Figure: active prime threshold ---
\begin{figure}[htbp]
\centering
\begin{tikzpicture}[xscale=2.0, yscale=4.5, >=stealth]
  % axes
  \draw[->,thick] (0,0) -- (6.5,0) node[right,font=\small] {prime $p$};
  \draw[->,thick] (0,0) -- (0,0.95) node[above,font=\small] {$\gamma_p$};
  % threshold line s = 1/2
  \draw[condorange, thick, dashed] (0,0.5) -- (6.3,0.5)
    node[right, font=\small] {$s = 1/2$};
  % bars
  \fill[thmblue!70] (0.7,0) rectangle (1.3,0.808);
  \node[above, font=\small, thmblue] at (1,0.808) {$0.808$};
  \node[below, font=\small] at (1,-0.02) {$3$};
  \fill[thmblue!70] (1.7,0) rectangle (2.3,0.696);
  \node[above, font=\small, thmblue] at (2,0.696) {$0.696$};
  \node[below, font=\small] at (2,-0.02) {$5$};
  \fill[thmblue!70] (2.7,0) rectangle (3.3,0.595);
  \node[above, font=\small, thmblue] at (3,0.595) {$0.595$};
  \node[below, font=\small] at (3,-0.02) {$7$};
  \fill[gray!50] (3.7,0) rectangle (4.3,0.426);
  \node[above, font=\small, gray] at (4,0.426) {$0.426$};
  \node[below, font=\small] at (4,-0.02) {$11$};
  \fill[gray!50] (4.7,0) rectangle (5.3,0.356);
  \node[above, font=\small, gray] at (5,0.356) {$0.356$};
  \node[below, font=\small] at (5,-0.02) {$13$};
  % labels
  \node[font=\small, thmblue] at (2,0.88) {\textbf{Active} ($\gamma_p > s$)};
  \node[font=\small, gray] at (4.5,0.88) {\textbf{Ghost} ($\gamma_p < s$)};
  % y-axis ticks
  \foreach \y in {0.2, 0.4, 0.6, 0.8} {
    \draw (-0.05, \y) -- (0.05, \y);
    \node[left, font=\scriptsize] at (-0.07, \y) {$\y$};
  }
\end{tikzpicture}
\caption{Anomalous dimensions $\gamma_p$ at $\mustar = 15$.
  The threshold $\gamma_p = s = 1/2$ (orange dashed) separates
  the three active primes $\{3,5,7\}$ (blue) from the ghost
  primes $\{11,13,\ldots\}$ (gray). This determines $N_c = 3$.}
\label{fig:active_threshold}
\end{figure}

% --- Figure: 3D coupling cube ---
\begin{figure}[htbp]
\centering
\tdplotsetmaincoords{70}{120}
\begin{tikzpicture}[tdplot_main_coords, scale=4.5]
  % Unit sphere (wireframe) — domain [0,1]^3 boundary
  % Equatorial circle in xy-plane (z=0.5)
  \draw[gray!25, thin] plot[domain=0:360, samples=72, smooth cycle]
    ({0.5 + 0.5*cos(\x)}, {0.5 + 0.5*sin(\x)}, 0.5);
  % Meridian circle in xz-plane (y=0.5)
  \draw[gray!25, thin] plot[domain=0:360, samples=72, smooth cycle]
    ({0.5 + 0.5*cos(\x)}, 0.5, {0.5 + 0.5*sin(\x)});
  % Meridian circle in yz-plane (x=0.5)
  \draw[gray!25, thin] plot[domain=0:360, samples=72, smooth cycle]
    (0.5, {0.5 + 0.5*cos(\x)}, {0.5 + 0.5*sin(\x)});
  % Axes
  \draw[->, thick, predred] (0,0,0) -- (1.25,0,0)
    node[below, font=\small] {$\sin^2\!\theta_3$};
  \draw[->, thick, thmblue] (0,0,0) -- (0,1.25,0)
    node[right, font=\small] {$\sin^2\!\theta_5$};
  \draw[->, thick, bridgepurple] (0,0,0) -- (0,0,1.25)
    node[above, font=\small] {$\sin^2\!\theta_7$};
  % The alpha_EM point: exact values from delta_p(2-delta_p) at q=13/15, mu*=15
  \pgfmathsetmacro{\sa}{0.2192}
  \pgfmathsetmacro{\sb}{0.1940}
  \pgfmathsetmacro{\sc}{0.1726}
  % Three circles intersecting at alpha point (radius 0.3)
  % Circle in sin^2(theta_3) = const plane (yz-plane at x=sa)
  \draw[predred, thick, fill=predred!8, fill opacity=0.3]
    plot[domain=0:360, samples=36, smooth cycle]
    (\sa, {\sb + 0.3*cos(\x)}, {\sc + 0.3*sin(\x)});
  % Circle in sin^2(theta_5) = const plane (xz-plane at y=sb)
  \draw[thmblue, thick, fill=thmblue!8, fill opacity=0.3]
    plot[domain=0:360, samples=36, smooth cycle]
    ({\sa + 0.3*cos(\x)}, \sb, {\sc + 0.3*sin(\x)});
  % Circle in sin^2(theta_7) = const plane (xy-plane at z=sc)
  \draw[bridgepurple, thick, fill=bridgepurple!8, fill opacity=0.3]
    plot[domain=0:360, samples=36, smooth cycle]
    ({\sa + 0.3*cos(\x)}, {\sb + 0.3*sin(\x)}, \sc);
  % Projection lines to axes
  \draw[predred, densely dashed, thin] (\sa,\sb,\sc) -- (\sa,0,0);
  \draw[thmblue, densely dashed, thin] (\sa,\sb,\sc) -- (0,\sb,0);
  \draw[bridgepurple, densely dashed, thin] (\sa,\sb,\sc) -- (0,0,\sc);
  % The coupling point
  \fill[condorange] (\sa,\sb,\sc) circle (1.0pt);
  % Leader line to label
  \draw[condorange, thin] (\sa,\sb,\sc) -- (0.55,0.55,0.65);
  \node[right, font=\small, condorange, align=left] at (0.55,0.55,0.65)
    {$\alphaEM = \prod \sinp{p}$\\[1pt]\footnotesize$= 1/137.036$};
  % Axis tick labels
  \node[below, font=\tiny, predred] at (\sa,0,0) {$0.219$};
  \node[left, font=\tiny, thmblue] at (0,\sb,0) {$0.194$};
  \node[left, font=\tiny, bridgepurple] at (0,0,\sc) {$0.172$};
  % Origin label
  \node[below left, font=\tiny, gray] at (0,0,0) {$0$};
\end{tikzpicture}
\caption{The coupling sphere.
  Each axis represents $\sinp{p}$ for one active prime.
  Three coloured circles, each perpendicular to one axis, intersect at the unique point
  $(\sinp{3}, \sinp{5}, \sinp{7})$ whose product equals
  $\alphaEM = 1/137.036$ (\BRIDGETAG, see~\cite{companion_M5}).
  The coupling point lies deep in the low-$\theta$ region:
  nature uses only a small fraction of the available parameter space.}
\label{fig:coupling_cube}
\end{figure}

% ============================================================
\subsection{$N_c = 3$: number of active primes}
\label{sec:Nc_3}

\mTHM
\begin{theorem}[$N_c = 3$]
\label{thm:Nc_3}
The number of active primes at $\mustar = 15$ is $N_c = 3$.

\begin{proof}
From Theorem~\ref{thm:active}: exactly three primes satisfy
$\gamma_p > s = 1/2$ at $\mustar = 15$, namely $\{3, 5, 7\}$.
All primes $p \geq 11$ have $\gamma_p < 1/2$.
Hence $N_c := |\{p : \gamma_p(\mustar) > s\}| = 3$.
\end{proof}
\end{theorem}

\begin{corollary}[Casimir coefficient]
\label{cor:CF}
$C_F = (N_c^2 - 1)/(2N_c) = 4/3$.
This is a \emph{consequence} of $N_c = 3$, not an input.
\end{corollary}

\begin{remark}[Algebraic consistency check (R14)]
\label{rem:Nc_algebraic}
The result can be cross-checked algebraically without assuming $C_F$.
T1 forbids 2 self-transitions at $p = 3$ (classes $1$ and $2$),
leaving $n_{\rm perm} = N_c^2 - 2 = 7$ permitted transitions out of~9.
Setting $n_{\rm perm}/N_c = C_F + 1$ with $C_F = (N_c^2 - 1)/(2N_c)$:
\[
  \frac{N_c^2 - 2}{N_c} = \frac{N_c^2 + 2N_c - 1}{2N_c},
\]
which simplifies to $N_c^2 - 2N_c - 3 = 0$, i.e.,
\begin{equation}
  (N_c - 3)(N_c + 1) = 0 \implies N_c = 3 \quad (N_c > 0).
  \label{eq:Nc_equation}
\end{equation}
Both the counting proof and the algebraic equation give $N_c = 3$
independently. (Verification: \texttt{ch06\_holonomy/test\_D07\_sin2\_identity.py}.)
\end{remark}

\begin{remark}[Five independent routes to $N_c = 3$]
\label{rem:Nc_routes}
Beyond the active-prime count and the algebraic check above,
three additional routes confirm $N_c = 3$ (R38b):
\begin{enumerate}[nosep]
\item $4N_c + 3 = \mustar$: Weyl fermion count per generation (BRIDGE).
\item $(N_c + 2)(N_c + 3) = 2\mustar = 30$: SU(5) decomposition
  $\bar{5} + 10 = 15$ (BRIDGE).
\item $(N_c + 1)!/(N_c + 3) = 2^{N_s - 1}$: only $N_c = 3$ gives
  an integer power of~2, yielding $N_s = 3$ spatial dimensions (BRIDGE).
\end{enumerate}
The active-prime count ($|\{p : \gamma_p > s\}| = 3$, \THMTAG) and
the algebraic check ($N_c = 3$ from the equation, \THMTAG) are
arithmetic results.  The three additional routes above (items 1--3)
involve physical identifications (\BRIDGETAG).
\end{remark}

\begin{remark}[Physical identification]
The identification of $N_c = 3$ (derived from sieve structure) with the
number of QCD colors is a BRIDGE claim (\cite{companion_M5}). Here we only prove
the arithmetic result that $N_c = 3$.
\end{remark}

% ============================================================
\subsection{Ghost primes and perturbative corrections}
\label{sec:ghost}

The inactive primes $\{11, 13, \ldots\}$ are not simply absent.
They appear as \emph{perturbative corrections} (ghost vacuum polarization).

\begin{definition}[Ghost VP weight]
\label{def:ghost_vp}
The ghost vacuum polarization weight is:
\begin{equation}
  \beta_{\rm ghost} = \sum_{p \in \{11,13\}} \sinsp{p}(\qrel) \cdot \gamma_p
  \approx 0.104.
  \label{eq:ghost_vp:ch06}
\end{equation}
\end{definition}

The ghost VP contribution to $\alpha_{\rm EM}$ is $\delta\alpha/\alpha
= -\gamma_3 \cdot \alpha_{\rm bare} \cdot \beta_{\rm ghost} \approx 14\;{\rm ppb}$
(see~\cite{companion_physics}).

% ============================================================
\subsection{The arithmetic origin of $\pi$}
\label{sec:pi_origin}

The holonomy identity reveals a structural fact about $\pi$:
the constant $\pi$ is not imported into the sieve from Euclidean
geometry. It \emph{emerges from} the arithmetic of divisibility.

\begin{remark}[Why the circle appears]
\label{rem:why_circle}
The transfer matrix $T_p$ acts on $\mathbb{Z}/p\mathbb{Z}$ by
redistributing probability mass among residue classes. Its $2 \times 2$
restriction to the non-zero classes (after T1 forces the diagonal
to zero at $p = 3$) has diagonal entries $\cos\theta_p = 1 - \delta_p$
and off-diagonal entries $\sin^2\!\theta_p = \delta_p(2 - \delta_p)$.
These satisfy the Pythagorean identity $\cos^2\theta + \sin^2\theta = 1$
\emph{automatically}: it is a consequence of the stochasticity constraint
$\sum_j T_{ij} = 1$ on the transfer matrix, not an assumption about
circles.

The unit circle $S^1$ is therefore the \emph{constraint manifold} of the
sieve's transition amplitudes. Its period $2\pi$ is determined by the
requirement that a full rotation returns to the identity. The sieve
does not use $\pi$---it \emph{generates} $S^1$, and $\pi$ is the
half-period of that circle.
\end{remark}

\begin{remark}[Computing $\pi$ from the sieve]
\label{rem:pi_computation}
The Euler product for $\zeta(2)$ admits a direct sieve interpretation.
At each prime~$p$, the sieve eliminates multiples of~$p$; the
\emph{squared} survival probability across all primes gives the
density of pairs coprime to every prime:
\begin{equation}
  \prod_{p\;\text{prime}} \frac{1}{1 - 1/p^2} = \zeta(2) = \frac{\pi^2}{6}.
  \label{eq:euler_pi}
\end{equation}
Each factor $(1 - 1/p^2)^{-1}$ is a sieve residue ratio: the
probability that two independent integers both survive the filter
at~$p$, divided by the unconstrained baseline. The infinite product
converges to $\pi^2/6$ --- a statement that the cumulative effect of
all arithmetic filters, composed multiplicatively, encodes the
geometry of $S^1$.

Truncating the product at the $k$-th prime $p_k$ gives a
\emph{sieve approximation} of $\pi$:
\begin{equation}
  \pi_k = \sqrt{6 \prod_{p \leq p_k} \frac{1}{1 - 1/p^2}},
  \qquad \pi_k \to \pi \;\;\text{as } k \to \infty.
  \label{eq:pi_sieve}
\end{equation}
Convergence is slow ($|\pi_k - \pi| \sim 3/p_k$), but the
conceptual point is sharp: $\pi$ is the limiting value of a purely
arithmetic product over primes, with no geometric input.
Verification: \texttt{ch09\_bridge/test\_equations\_pont\_PT.py}, test~H5 (24
primes, $|\pi_{24} - \pi|/\pi < 0.5\%$).
\end{remark}

\begin{remark}[Why $\pi$ appears in prime number formulas]
\label{rem:why_pi_in_primes}
The classical surprise---why does $\pi$ appear in $\zeta(2)$, in
Mertens' theorem, in the Hardy--Littlewood singular series, in the
prime number theorem error term?---dissolves in the holonomy framework.

At every prime $p$, the sieve imposes a \emph{rotation} $\theta_p$
on the residue space. The cumulative holonomy of these rotations
is measured by products of trigonometric functions of $\theta_p$.
Since $\sin$ and $\cos$ are periodic with period $2\pi$, every
multiplicative formula over primes inherits~$\pi$ as the
\emph{period of the rotation that divisibility generates}.

In particular:
\begin{itemize}[nosep]
\item $\zeta(2) = \pi^2/6$: cumulative squared-survival product
  (\cref{eq:euler_pi}).
\item $\prod_{p \leq x}(1 - 1/p) \sim 2e^{-\gamma}/\ln x$
  (Mertens): the survival fraction decays logarithmically, and
  the constant $e^{-\gamma}$ encodes the same holonomy accumulation.
\item $\alphaEM = \prod_{p \in \{3,5,7\}} \sinp{p}$: the
  fine-structure constant is a product of $\sin^2\!\theta_p$ values
  --- literally a product of holonomy amplitudes on~$S^1$.
\item $\mu_{\rm end} = N_c \cdot \pi = 3\pi$: the RG integration
  bound equals $N_c$ half-turns (Berry phase, see~\cite{companion_physics}).
\end{itemize}

In every case, $\pi$ enters number theory not as an import from
geometry, but as the \emph{intrinsic period of the rotational
dynamics that the sieve creates on residue spaces}. The
arithmetic constrains the geometry; the geometry returns $\pi$.
\end{remark}

% ============================================================
\subsection{Summary and connections}
\label{sec:ch6_summary}

\begin{ledgerbox}
Hard core:
  Holonomy Identity: $\sinp{p} = \delta_p(2-\delta_p)$ -- algebraic identity.
  Active prime criterion: $\{3,5,7\}$ active ($\gamma_p > s$) -- numerical THM at $\mustar=15$.
  $N_c = 3$: counting active primes at $\mustar = 15$ -- arithmetic THM.
  Algebraic check: $(N_c-3)(N_c+1)=0$ from T1 transition count.

Conditional:
  Numerical values of $\sinp{p}$ and $\gamma_p$ use $\mustar = 15$ (\cite{companion_M3}, T5 by scan).

Identification (THM, four unicities):
  $\sinp{p}$ = coupling constant (\cite{companion_M5}, BA3, closed by G1+G3).
  $N_c = 3$ = QCD colors (\cite{companion_M5}, closed by T5+T6).
  $\{11,13\}$ = ghost VP counter-terms (\cite{companion_physics}, DER-PHYS).

Validation:
  $\sinp{p}$ verified by \texttt{ch06\_holonomy/test\_D07\_sin2\_identity.py}.
  $N_c = 3$ verified by \texttt{ch06\_holonomy/test\_D07\_sin2\_identity.py}.
  Ghost VP 14\,ppb verified (\cite{companion_physics}).
\end{ledgerbox}

\medskip\noindent\textbf{Forward connections.}
\begin{itemize}
\item \cite{companion_M3} provides $\mustar = 15$ (needed for numerical values here).
\item \cite{companion_M5} derives $\sinsp{p}$ as coupling constants (DER, Pontryagin + Wightman).
\item The companion article~\cite{companion_physics} assembles the product $\prod_{p=3,5,7} \sinsp{p}$
  to derive $\alpha_{\rm EM}$.
\item The companion article~\cite{companion_physics} uses $\gamma_p$ to derive $\sinW$ and $\alpha_s$.
\item The companion article~\cite{companion_physics} uses the active directions to define the spatial metric.
\end{itemize}


\newpage


% ============================================================
%  CONCLUSION
% ============================================================
\section{Conclusion}
\label{sec:conclusion}

This article has established two pillars of the Theory of Persistence
that connect the sieve's combinatorial structure to geometric and
spectral objects.

\textbf{Information geometry.}
The CRT product structure of the sieve state space $\Delta_m$ uniquely
selects two geometric objects: the KL divergence $\DKL$ (Theorem~G1,
via Shore--Johnson reduction and an autonomous CRT proof) and the
Fisher metric $\Fisher$ (Theorem~G3, via \v{C}encov reduction).  The
connection between them is standard (Theorem~G2: Fisher is the Hessian
of $\DKL$).  The key PT contribution is not in proving Shore--Johnson
or \v{C}encov---these are imported classical theorems---but in
verifying that the sieve state space, with its CRT factorization and
Markov transfer matrices, satisfies the premises of both uniqueness
theorems.  Seven elimination tests confirm that no alternative
divergence or metric is compatible with the sieve's symmetries.

The Fisher metric encodes the duality between addition and
multiplication: it is the unique geometry that respects both
operations simultaneously.  The GFT identity $\log_2 m = \DKL + H$
is this duality expressed as a conservation law; Fisher geometry is
its incarnation as a Riemannian metric.

\textbf{Holonomy.}
The holonomy identity $\sinsp{p} = \delta_p(2 - \delta_p)$ is an
algebraic tautology that receives three independent derivations:
geometric (Pythagorean identity from stochasticity), spectral (Fourier
contraction of the fundamental character), and information-geometric
(per-prime Fisher curvature component).  The agreement of all three
routes eliminates single-proof risk.

The anomalous dimensions $\gamma_p$ are strictly decreasing in $p$,
and the active prime criterion $\gamma_p > s = 1/2$ selects exactly
$\{3, 5, 7\}$ at the fixed point $\mustar = 15$
(see~\cite{companion_M3}).  This yields $N_c = 3$ as a theorem, not
an input.  The Casimir coefficient $C_F = 4/3$ follows as a
corollary.  An independent algebraic check---$(N_c - 3)(N_c + 1)
= 0$ from the T1 transition count---confirms the result without
reference to anomalous dimensions.

The three active primes define three independent holonomy directions.
Their physical identification as spatial dimensions, and the
identification of $N_c = 3$ with the number of QCD colours, are
bridge claims developed in~\cite{companion_M5}.

\textbf{Open questions.}
The threshold criterion $\gamma_p > s$ is natural (five independent
arguments, 35\%-wide stability interval) but lacks a closed-form
optimality proof.  A rigorous derivation showing that $\tau = s$ is
the \emph{only} consistent threshold---rather than merely the most
natural choice within a wide stability band---remains open.

The ghost primes ($p \geq 11$) contribute perturbatively as vacuum
polarization corrections.  Whether the ghost VP series converges
absolutely, and whether it admits a closed-form resummation, are
questions for future work.

\textbf{Forward.}
The companion article~\cite{companion_M3} proves convergence of the
sieve recurrences and establishes the fixed point $\mustar = 15$ that
underpins the numerical values used throughout this article.
The extended mathematical structures---complex mechanics, spectral
graphs, and symmetry extensions---are developed
in~\cite{companion_M4}.  The bridge from arithmetic to physics, where
the geometric objects established here acquire physical meaning, is
the subject of~\cite{companion_M5}.

\bibliographystyle{plainnat}
\bibliography{references}

\end{document}
